% SGA 3, 第 VI_B 篇, 第 1 轮英文重构.
% 范围: 例 5.8.3 全文.
% 底本: Polo--Gille Exp6B-13oct24.pdf, 分篇第 35--36 页
% (合订本第 376--377 页; 原版页边印刷页码 357, 重版本页眉 365--366).
% 例 5.8.4 始于分篇第 36 页 / 合订本第 377 页.
% OCR 与 Reinhold 英文 Markdown 仅为比对材料.
% 本单元无图及底本编者注.

\medskip
\phantomsection\label{ex:sga3-VIB-5.8.3}%
\noindent\textbf{例 5.8.3.---}
固定一个代数闭域 \(k\). 先回忆, 每个``局部 Boolean''拓扑空间 \(X\),
即每个有紧开集基的空间, 配备局部常值的 \(k\)-值函数层后,
就成为一个 \(k\)-概形
(参见 [DG70], I, \S1, \SGARefLink{cor:sga3-VIB-2.12}{2.12}).
这样的 \(X\) 称为\emph{局部 Boolean \(k\)-概形}.

还注意到, 每个有分离开集基的拓扑空间 \(E\),
以及每个概形 \(X\), 都有一个稠密分离开集.
事实上, 分离开集的每个递增并都是分离的
(对概形 \(X\), 这由 \hyperref[lem:sga3-VIB-5.6.3]{5.6.3} 得到),
因此存在一个极大的这样的开集 \(U\).
它必稠密: 否则就有非空开集 \(V\), 使 \(U\cap V=\varnothing\);
开集 \(V\) 将包含一个非空分离开集 \(W\), 而
\[
  U\cup W
\]
仍然分离, 与 \(U\) 的极大性矛盾.

现设 \(C\) 是 Cantor 空间, 看成赋予乘积拓扑的群
\[
  (\mathbb Z/2\mathbb Z)^{\mathbb N}
\]
的底空间. 对每个 \(p\in C\), 设 \(C(p)\) 是 \(C\) 的另一个副本,
并设 \(X\) 由将每个 \(C(p)\) 沿 \(C-\{p\}\) 与 \(C\) 粘合得到.
则 \(X\) 是不分离的局部 Boolean \(k\)-概形,
且 \(C\) 是 \(X\) 的稠密开集.

设 \(G\) 是 \(k\)-概形 \(X\) 的自同构群, 即 \(X\) 的同胚群.
则 \(G\) 在 \(X\) 上传递地作用.
事实上, \(X\) 是 \(C\) 与每个点 \(p\in C\) 的第二个点 \(p'\) 的并.
这个第二点的特征是: 它是使
\[
  (C-\{p\})\cup\{x\}
\]
分离的唯一的 \(x\in X-\{p\}\).
因此 \(C\) 的每个自同构 \(\phi\) 都延拓为 \(X\) 的一个自同构
\(\phi_X\), 它对每个 \(p\) 满足
\[
  \phi_X(p')=\phi_X(p)'
\]
此外, 交换 \(p\) 与 \(p'\) 并固定其他所有点的映射
\[
  \theta_p:X\longrightarrow X
\]
是 \(X\) 的自同构. 最后, \(C\) 的自同构群在 \(C\) 上传递地作用:
利用 \(C\) 的群结构, 只须考虑平移即可.

因此, 局部有限型的离散 \(k\)-群 \(G\) 作用于 \(k\)-概形 \(X\), 使
\[
  G\times X\longrightarrow X\times X,
  \qquad (g,x)\longmapsto(gx,x),
\]
满射, 但 \(X\) 不分离, 尽管 \(C\) 是一个稠密分离开集.
